%----------------------------------------------------------------
% Poglavje (Chapter) 2
%----------------------------------------------------------------
\chapter{Related work}
\label{ch:related_work}

The research into the relationship between data serialization and performance of LLMs has grown from a prompt
engineering concern, to the important questions a lot of experts are asking. As LLMs are increasingly
integrated into enterprise pipelines, the way data is structured, be it JSON, XML, CSV, Markdown, has been
shown to hold a significant influence over the overall performance of the models (cost, accuracy, efficiency).
This is not a mere question of our format preference, the benchmarks such as StrucText-Eval
\cite{struct-text-eval} has given us insight how such formats influence the overall accuracy of the models.

\section{A Survey of Structured Data Representations}

To understand how format influences models, it's paramount to have a way of organizing structured data
representations that we are going to be using. We can identify several main categories of data representation, as
shown by the StrucText-Eval \cite{struct-text-eval}. These categories show the range of data formats, from
those that are strictly structured to those that are semi-structured. Structured data is consistent and
organized. This makes it easy for computers
to process, but it also creates some overhead when LLMs are tokenizing the data. The most common
classification system divides these into four main types:
% TODO: Add citations for all of the data formats !!!
\begin{itemize}
  \item \textbf{Tree and Tabular Formats:} These includes CVS files and custom hierarchical notation. In these
    formats, relationships are defined either by positional alignment (as in CSV) or explicit pointers (as in
    tree structures). However, CSV fundamentally cannot represent nested structures, limiting its
    applicability to flat tabular data.
  \item \textbf{Object Notation:} This category represents the most common data interchange languages in
    software development, notably JSON (JavaScript Object Notation), YAML (YAML Ain't Markup Language), and
    TOML. These formats rely on key-value pairs and nesting to represent multi-dimensional data
  \item \textbf{Markup Languages:} Semi-structured data types like Markdown, LaTeX, and Org represent
    information in a human-readable form while embedding structural metadata (headings, lists, and semantic
    tags). Markdown has become the de facto standard for general-purpose LLM interactions due to its presence
    in training data.
  \item \textbf{Code Based Formats:} These include HTML, JSX, SVG, and XML. These formats are special
    because they are measured not only by how well they follow the rules of language, but also by how well they
    can create a visual or logical output that meets the requirements of the user.

\end{itemize}

% The categorization of these formats is critical because each imposes a different "tokenization
% overhead." For instance, JSON requires significant structural punctuation (braces, quotes, colons), whereas
% YAML relies on whitespace and indentation. The structural density of a format influences the model's ability
% to maintain context over long sequences, as redundant tokens can "dilute" the attention allocated to the
% actual semantic content.

\section{Token Efficiency and Economic Impact of Data Serialization}

The choice of data format has direct impact on the computational cost of LLMs. Format verbosity determines
token token usage of the model, and this has a knock on effect for other factors, such as latency, API cost,
and effective context window.

Formats like JSON and XML are inefficient with token usage, as they have been designed for data exchange
between information systems, where self description and compatibility were more important than their
verbosity. These formats are designed in such a way that each fields carries its own description, through
repeated key words, which create overhead when working with LLMs. The experiments performed in \cite{}, were
carried out across $1.000$ IoT telemetry records. They have shown that majority of encoding tokens used for
encoding data in the JSON format is used up by field names and punctuation ($60.500$ tokens), while the
actual semantic values take up only $14.274$ tokens, out of $79.774$ tokens. This overhead is problematic in
the current age, as most pricing is done based on the input and output tokens used, with this higher token
usage directly translates to higher operational costs for the users. Also in the same experiment it was shown
that reduction in tokens used directly correlates to lower latency, when compared to their custom ONTO format.

In order to address these inefficiencies, there have been multiple efforts aimed at them, in form of ONTO
(Object Notation for Token Optimization) and TOON (Token Optimized Object Notation). Both formats are
oriented at fixing the problem of repeated schema for each record, they fix this by declaring all the fields
only once at the top of the record, and values are arrange in pipe delimited rows with indentation based
hierarchy. Both of these models claim to use significantly less tokens (TOON $39.9\%$ fewer tokens compared
to JSON, and ONTO $46-51\%$ compared to JSON ) while minimal impacts on accuracy.
